\writeSectionFrame{\presentationTitle}{Das Problem}
\begin{frame}{Iterieren durch Mengen}
	\begin{center}
 		\includegraphics[scale=0.4]{\BILDERPFAD/problem1.png}
 	\end{center}
\end{frame}

\note{
 	Das Navigieren durch Elemente ist eines der ältesten Probleme der Informatik.
 	Gelöst wurde es immer mit Arrays und Schleifen, wobei der Programmierer dann
 	auf Dinge achten musste wie Grenzen u.ä.
}

\begin{frame}[fragile]{Lösung in modernen Programmiersprachen}
	\begin{exampleblock}{StandardIteratorBeispiel.java}
        \begin{lstlisting}
List<String> zauberspruch = new ArrayList<>();
zauberspruch.add("Abra");
zauberspruch.add("kada");
zauberspruch.add("bra");

for (String myString: zauberspruch)  {
    System.out.print(myString);
}
        \end{lstlisting}
 	\end{exampleblock}
\end{frame}

\note{
	Moderne Programmiersprachen bieten den Iterator als einheitliche Schnittstelle
	für unterschiedliche Datenstrukturen an. Die for-each-Schleife funktioniert für 
	alle Klassen, die Iterable implementieren. 
}

\frameWithImageOnly{\BILDERPFAD/problem3.jpg}
\note{
	Da stellt sich die Frage: Warum sollte man dann selbst das Muster noch implementieren, wenn
	die Programmiersprache bereits eine Implementierung mitbringt?
}

\begin{frame}{Iterieren durch komplexe Datenstrukturen}
	\begin{center}
 		\includegraphics[scale=0.25]{\BILDERPFAD/problem4.png}
 	\end{center}
\end{frame}

\note{
 	Man kann auch durch komplexe Datenstrukturen iterieren wollen und nicht nur durch
 	flache Listen. Damit hätten wir eine natürliche Kombinierbarkeit mit dem 
 	Kompositum-Muster.
}

\begin{frame}{Logik des Iterierens}
	\begin{center}
 		\includegraphics[scale=0.5]{\BILDERPFAD/problem5.png}
 	\end{center}
\end{frame}

\note{
 	Außerdem kann ein Iterator auch Geschäftslogik enthalten, so dass die Logik für die
 	Auswahl bestimmter Elemente im Iterator gekapselt werden kann.
}

\frameWithImageOnly{\BILDERPFAD/problem6.jpg}
\note{
	Vorneweg kann man schon einmal sagen, dass das Iterator-Muster eines der flexibelsten und
	am besten anpassbaren ist. 
}


\frameWithImageOnly{\BILDERPFAD/problem7.jpg}
\note{
	Stellen wir uns vor, wir machen Urlaub in Rom. Wir können nun auf eigene Faust losziehen und auf der
	Suche nach dem Kollosseum Kreise und Umwege laufen. Wir könnten eine Navigationsapp auf unserem Handy
	benutzen, aber wir könnten auch einen Fremdenführer buchen. Alle Methoden haben ihre Vor- und Nachteile
	und letztendlich entspricht es unserer persönlichen Vorliebe, welche Methode wir verwenden. Was alle drei
	Möglichkeiten gemeinsam haben: Sie alle ermöglichen es uns, durch eine Menge von Sehenswürdigkeiten zu laufen
	und letztendlich ist für uns die genaue Implementierung dieses Durchlaufens erst einmal zweitrangig. Wir 
	möchten nur irgendwie das Kollosseum, den Petersdom, die Engelsburg usw. sehen. Diese drei Möglichkeiten
	sind verschiedene Iteratoren durch die Menge von Sehenswürdigkeiten.
}

